\chapter{Proofs for Huber's General Information Game}
\label{annex:p15-huber-information-game}

\begin{annexitem}{annex:p15.vague.subprobability.compactness}{thm:p15-vague-subprobability-compactness}{The compact domain and the closure of contamination}
\begin{proof}
Choose a countable collection \(h_j\) of compactly supported continuous
piecewise linear functions with rational vertices and values. Include
enough such functions that every \(h\in C_c(\mathbb R)\) has uniform
approximants from the collection with supports in one bounded interval.
Linear interpolation and rational approximation construct these
approximants. Put
\[
 d(F,G)=\sum_{j\geq1}2^{-j}
       \min\left(1,\left|\int h_j\,dF-\int h_j\,dG\right|\right).
\]
If all these integrals agree, uniform approximation gives equality
for every \(C_c\) function. Approximate the indicator of a bounded
interval whose endpoints are not atoms of \(F+G\) by continuous ramps.
Dominated convergence gives equality on those intervals. There are
only countably many atoms, so such endpoints form a dense set; these
intervals determine a finite Borel measure by monotone limits and
the uniqueness of extension from the interval algebra. Thus \(d\)
is a metric. Convergence for \(d\) implies vague convergence because
the error in either measure's integral of a uniform approximation
is at most the uniform error. Vague convergence implies convergence
for \(d\) by first controlling finitely many coordinates and then
the geometric tail.

For a sequence \(F_n\in\mathcal M_{\leq1}\), make proper probabilities
\(P_n=F_n+(1-F_n(\mathbb R))\delta_n\). The Helly construction in
\cref{thm:p15-helly-selection} gives a subsequence and a nondecreasing
right-continuous function with limits \(\beta\leq\gamma\), inducing a
measure \(F\) of mass \(\gamma-\beta\leq1\). The Stieltjes construction
is \cref{thm:p15-stieltjes-construction}. Intervals with endpoints
at continuity points have convergent masses along this subsequence.
Approximate a \(C_c\) function on a bounded interval by step functions
on a partition with such endpoints. Its error is bounded by its
modulus of continuity times a mass at most one. This proves
\(\int h\,dP_n\to\int h\,dF\). For fixed compact support the added
atom is eventually outside it, giving \(F_n\to F\) vaguely.
The metric space is therefore sequentially compact. For completeness
of the topological step, a sequentially compact metric space is totally
bounded, since an infinite sequence of points separated by a fixed
positive distance otherwise contradicts sequential compactness.
It is complete, since a convergent subsequence of a Cauchy sequence
forces that whole sequence to converge. A complete totally bounded
space is compact: successive finite covers by balls of radii tending
to zero either give a finite subcover of a proposed open cover or
produce a Cauchy sequence whose limit contradicts the absence of one.

Write \(a=1-\varepsilon\). A vague limit of measures dominating \(aG\)
satisfies \(\int h\,dF\geq a\int h\,dG\) for every nonnegative \(h\in C_c\).
This implies \(F\geq aG\): continuous ramps first give the inequality
on open bounded intervals and their finite unions, and increasing
unions and finite-measure approximation extend it to Borel sets.
One may justify that last approximation directly: the class of sets
approximable in the measure \(F+aG\) by the interval algebra is closed
under complements and countable unions, by continuity of finite
measures and finite-union approximation. It contains the generating
intervals and hence all Borel sets.
Conversely such an \(F\) is \(aG+\varepsilon H\) with a nonnegative
measure \(H\) of mass at most one. Replace \(H\) by
\(H+(1-H(\mathbb R))\delta_n\) to approximate it by probabilities.
This proves the claimed closure. Convexity is immediate and the
domination inequalities define a closed set. Symmetry is the closed
condition \(\int h(x)\,dF=\int h(-x)\,dF\); replacing the escaping atom
by \((\delta_n+\delta_{-n})/2\) proves the symmetric version.
The proper class is not closed: take \(H=\delta_n\).
\end{proof}
\end{annexitem}

\begin{annexitem}{annex:p15.information.payoff.bridge}{thm:p15-information-payoff-bridge}{The original ratio, its affine representation and mixtures}
\begin{proof}
For fixed \(\psi,F\), abbreviate \(A=A_\psi(F)\) and \(B=B_\psi(F)\).
The supremum over real \(s\) of \(2sA-s^2B\) is \(A^2/B\) if \(B>0\),
infinity if \(B=0,A\ne0\), and zero otherwise. Since \(s\psi\in\mathcal D\),
taking suprema proves \(I=\sup J\). If \(F\ne0\), some
\(\eta\in\mathcal D\) has \(B_\eta>0\), by a bump positive on an interval
of positive mass. When \(B_\psi=0,A_\psi\ne0\), \(\psi=0\) \(F\)-almost
everywhere, and
\[
 B_{\psi+t\eta}=t^2B_\eta,\qquad
 A_{\psi+t\eta}=A_\psi+tA_\eta .
\]
Their ratios tend to infinity as \(t\to0\). Thus restricting the
supremum to positive denominators loses nothing for nonzero \(F\).
Each \(J(\psi,\cdot)\) is continuous and affine in the vague topology.
Its supremum is lower semicontinuous and convex.

For the displayed perspective inequality put
\(v=tv_1+(1-t)v_2\) and \(w=tv_1/v\).
The left side is
\(v(wu_1/v_1+(1-w)u_2/v_2)^2\); convexity of a square bounds this
by the asserted right side. Zero-denominator cases follow either
by limits when the numerator is zero or directly from
\(A^2/B=\sup_s(2sA-s^2B)\) with the stated conventions.
This also proves convexity of \(K\).

Here are the measure details for an arbitrary mixed strategy.
For open \(O\), continuous compactly supported functions increasing
to \({\bf1}_O\) show that \(F\mapsto F(O)\) is Borel measurable.
Complements and increasing disjoint unions extend measurability to
all Borel sets. Therefore \(\overline F\) is defined; monotone
convergence gives countable additivity and mass at most one.
Partition the compact metric space \(C\) into finitely many Borel
sets of diameter at most \(1/n\), and choose representatives.
The corresponding finite convex combinations \(F_n\) belong to \(C\).
For each \(h\in C_c\), the continuous function
\(F\mapsto\int h\,dF\) is uniformly continuous on \(C\).
It follows that \(F_n\to\overline F\) vaguely, so \(\overline F\in C\).
Finally, if \(\int K\,d\Pi<\infty\), then \(A=0\) on the set \(B=0\),
up to a \(\Pi\)-null set, and Cauchy--Schwarz gives
\[
 \left(\int A\,d\Pi\right)^2
 \leq \left(\int_{\{B>0\}}A^2/B\,d\Pi\right)\left(\int B\,d\Pi\right).
\]
If \(\int B=0\), both averaged numerator and denominator are zero.
If \(\int K=\infty\), the desired inequality is automatic.
These cases establish the barycenter inequality in full.
\end{proof}
\end{annexitem}

\begin{annexitem}{annex:p15.finite.information.characterization}{thm:p15-finite-information-characterization}{All of the finite-information regularity conclusion}
\begin{proof}
We first supply the functional-analysis ingredients in the particular
finite-measure setting needed here. The functions \(C_c^1\) are dense
in \(L^2(F)\). To see this, choose a dense set of interval endpoints
avoiding the countably many atoms of \(F\). Finite unions of intervals
with those endpoints approximate every Borel set in \(F\)-measure:
the class of sets so approximable is closed under complements and
countable unions, the latter by first truncating to a finite union.
It contains an algebra generating the Borel sets. Tail truncation
reduces to bounded intervals. Smooth ramps converging to their
indicators converge in \(L^2(F)\), since the endpoints are not atoms.
Approximating simple functions and then truncating arbitrary \(L^2\)
functions proves the claim.

We will also use completeness of \(L^2(F)\), with the following direct
proof. From a Cauchy sequence select a subsequence whose successive
differences have summable \(L^2\) norms. The sums of their absolute
values have bounded \(L^2\) norm by the triangle inequality.
Monotone convergence shows that the infinite sum is finite almost
everywhere and in \(L^2\); its tails tend to zero in \(L^2\).
Thus the subsequence, and consequently the original sequence,
has an \(L^2\) limit.

Suppose \(I(F)<\infty\). The functional
\(L(\psi)=\int\psi'\,dF\) satisfies
\[
 |L(\psi)|\leq\sqrt{I(F)}\,\|\psi\|_{L^2(F)}.
\]
It is well defined on \(L^2\) equivalence classes of test functions:
zero norm forces \(L=0\). It therefore extends by the preceding
density to a bounded linear functional on \(L^2(F)\).
We give its representation without presupposing a representation
theorem. If \(L=0\), take \(u=0\). Otherwise minimize the squared norm
on the closed affine set \(L(v)=1\). Its infimum is positive by
boundedness of \(L\). The parallelogram identity applied to two members
of a minimizing sequence and their midpoint makes that sequence
Cauchy; completeness supplies a minimizer \(v_0\).
Varying \(v_0+t w\), \(w\in\ker L\), shows
\(\langle v_0,w\rangle=0\). Hence
\[
 L(v)=\langle v,u\rangle,\qquad
 u=v_0/\|v_0\|_2^2 .
\]
The operator norm squared, equivalently the supremum of the test
ratios by density, is \(\|u\|_2^2=I(F)\).

Let \(\nu=-uF\). It is a finite signed measure, since
\(\int|u|\,dF\leq\sqrt{F(\mathbb R)I(F)}\).
Put \(H(x)=\nu((-\infty,x])\). Fubini, applied to the positive
and negative variations of this finite measure, gives
\[
 \int\psi'(x)H(x)\,dx=-\int\psi\,d\nu
                  =\int\psi'\,dF\qquad(\psi\in C_c^1).
\]
Thus the locally finite signed measure \(T=F-H\,dx\) annihilates
derivatives of \(C_c^1\) functions. If \(\varphi\in C_c\) has
Lebesgue integral zero, its primitive
\(\psi(x)=\int_{-\infty}^x\varphi(t)\,dt\) belongs to \(C_c^1\);
therefore \(\int\varphi\,dT=0\).
Fix \(\eta\in C_c\) of integral one and put \(c=\int\eta\,dT\).
Subtracting \(\eta\int\varphi\,dx\) shows that
\(\int\varphi\,dT=c\int\varphi\,dx\) for every \(\varphi\in C_c\).
Continuous compactly supported functions determine locally finite
measures, by the same interval-ramp argument on each bounded interval
used above. Consequently \(F=(H+c)\,dx\).

Let \(f=H+c\) initially almost everywhere. Since now
\(\nu=-uf\,dx\), its cumulative function \(H\) is absolutely continuous,
with derivative \(-uf\in L^1(dx)\). Choose this continuous representative
of \(f\). It is nonnegative everywhere, since it is nonnegative almost
everywhere and continuous. We have
\[
 f'=-uf\quad\hbox{a.e.},\qquad
 \int_{\{f>0\}}f'^2/f=\int u^2\,dF=I(F).
\]
This proves absolute continuity as a function, rather than merely
absolute continuity of the measure.
The derivative bound follows by Cauchy--Schwarz. Its integrability
gives finite limits of \(f\) at both infinities; nonnegativity and
integrability of \(f\) force both limits to be zero. Total variation
is at least \(2f(x)\) for each \(x\). This proves the bound on
\(\|f\|_\infty\), continuity and bounded variation. If \(I(F)=0\),
then \(f'=0\) almost everywhere, so the integrable constant \(f\)
is zero.

Conversely let \(f\) have the asserted properties and finite classical
information. A nonnegative absolutely continuous function has
derivative zero almost everywhere on its zero set: at every
differentiability point there, it has a local minimum.
Set \(s=-f'/f\) on \(\{f>0\}\) and zero elsewhere.
Integration by parts against a compactly supported \(C^1\) function
and Cauchy--Schwarz give
\[
 \int\psi' f=\int\psi s f,\qquad
 K(\psi,F)\leq\int s^2 f .
\]
Approximate \(s\) in \(L^2(F)\) by test functions \(\psi_n\).
Then \(\int\psi_n s f\to\int s^2f\) and
\(\int\psi_n^2f\to\int s^2f\). Their ratios converge to the latter
integral when it is positive; the zero case has already been handled.
Thus the variational and classical information coincide.
These arguments supply the density, boundedness, continuity,
absolute continuity, bounded variation and information conclusions
used in the source's Theorem 3.
\end{proof}
\end{annexitem}

\begin{annexitem}{annex:p15.information.game.value}{thm:p15-information-game-value}{Existence of minimizers and the exact minimax value}
\begin{proof}
Lower semicontinuity makes the sets
\(\{F\in C:I(F)\leq\alpha+1/n\}\) closed, nested and nonempty.
Compactness supplies a point in their intersection, an information
minimizer \(F_*\).
Fix \(r<\alpha\). At each \(F\in C\), some test function has
\(J(\psi,F)>r\). Continuity and compactness give a finite subcover
using \(\psi_1,\ldots,\psi_m\).
The compact convex set
\[
 D=\{(J(\psi_1,F),\ldots,J(\psi_m,F)):F\in C\}
\]
is disjoint from \(O=(-\infty,r]^m\). Their distance is positive.
Choose a nearest pair \(d_0\in D,o_0\in O\). The nearest point of
\(O\) to \(d_0\) has coordinates \(\min(d_{0j},r)\), so
\(v=d_0-o_0\) is nonnegative and nonzero.
Varying \(d_0\) toward any \(d\in D\) gives
\(v\cdot d\geq v\cdot d_0=r\sum v_j+\|v\|^2\).
With \(w_j=v_j/\sum v_j\) and \(\psi=\sum w_j\psi_j\), convexity
of the square gives
\[
 J(\psi,F)\geq\sum_jw_jJ(\psi_j,F)>r\qquad(F\in C).
\]
This elementary finite-dimensional separation proves the required
minimax inequality without assuming an infinite-dimensional minimax
theorem. The reverse inequality follows by evaluating at \(F_*\).
Let \(r\uparrow\alpha\).

For every score \(K\geq J\): when \(B>0\), their difference is
\((A-B)^2/B\); the zero cases follow directly. Conversely
\(\inf_C K(\psi,\cdot)\leq K(\psi,F_*)\leq I(F_*)=\alpha\).
This proves the ratio-game equality, including \(\alpha=0\).
For a standard normal density \(\phi\), integration by parts gives
\(A_\psi=\int x\psi(x)\phi(x)\,dx\).
Cauchy--Schwarz bounds its ratio by \(\int x^2\phi=1\), with equality
only when \(\psi=cx\) almost everywhere for a nonzero \(c\).
No such function has compact support. Approximating \(x\) in
\(L^2(\phi\,dx)\) by test functions gives the value one.
\end{proof}
\end{annexitem}

\begin{annexitem}{annex:p15.information.equality.components}{thm:p15-information-equality-components}{The exact equality condition replacing blanket uniqueness}
\begin{proof}
At points where \(f,g>0\), direct common-denominator expansion gives
\[
 t\frac{f'^2}{f}+(1-t)\frac{g'^2}{g}
 -\frac{(tf'+(1-t)g')^2}{tf+(1-t)g}
 =\frac{t(1-t)fg}{tf+(1-t)g}
       \left(\frac{f'}f-\frac{g'}g\right)^2 .
\]
At almost every zero of a nonnegative absolutely continuous function
its derivative is zero. Thus the deficit vanishes off the overlap,
with the usual zero-density convention. Integrating proves the identity
and its almost-everywhere equality criterion.

For the componentwise assertion put \(w=f+g\).
On a compact subinterval of a connected component of \(\{w>0\}\),
continuity bounds \(w\) away from zero, so \(f/w\) is absolutely
continuous there. Its derivative is
\((f'g-fg')/w^2\). Under equality this is zero on the overlap and
also almost everywhere where one function vanishes.
It follows that \(f/w\) is constant throughout each component.
The converse follows from the same formula.

If \(f,g\) are minimizers with value \(\alpha\), their mixture has
information both at least and at most \(\alpha\); hence equality holds.
If the union of positive supports is connected, write
\(f=a w,g=(1-a)w\) there. Information is homogeneous, so
\(\alpha=aI(w)=(1-a)I(w)\). When \(\alpha>0\), this forces
\(a=1/2\) and \(f=g\). The characterization of zero information
handles that remaining case. If one minimizer is positive everywhere,
the union is the whole line, proving its uniqueness.
\end{proof}
\end{annexitem}

\begin{annexitem}{annex:p15.information.first.variation}{thm:p15-information-first-variation}{The support term in the source's variational condition}
\begin{proof}
Use zero for \(f'^2/f\) where \(f=0\), and similarly for other
nonnegative absolutely continuous densities. Let \(i_t\) be the
pointwise information integrand of \((1-t)f+tg\).
Perspective convexity makes
\(d_t=(i_t-i_0)/t\) nondecreasing as \(t\) increases and bounds it
above by \(i_1-i_0\), an integrable function.
For \(t\downarrow0\), monotone convergence applied to the nonnegative
functions \(i_1-i_0-d_t\) permits integration of the pointwise limit,
possibly with value minus infinity.
On \(S=\{f>0\}\) differentiation gives
\[
 d_0=-2\psi g'-\psi^2g-f'^2/f ,
\]
whereas on \(S^c\), \(f=f'=0\) almost everywhere and \(d_0=g'^2/g\).
There is therefore an actual contribution from newly occupied support.

If \(b_g<\infty\), Cauchy--Schwarz gives
\[
 \int_S|\psi g'|\leq
 \left(\int_S\psi^2g\right)^{1/2}
 \left(\int_{\{g>0\}}g'^2/g\right)^{1/2}<\infty,
\]
and integration yields \(2a_g-b_g+r_g-I(f)\).
If \(b_g=\infty\), the inequality
\(-2\psi g'\leq\frac12\psi^2g+2g'^2/g\), valid also at almost
every zero of \(g\), bounds
\(d_0\leq-\frac12\psi^2g+2g'^2/g-f'^2/f\) on \(S\).
Its positive part is integrably bounded, while its integral is minus
infinity. This establishes the claimed extended derivative.

A minimum of a convex function along each segment has a nonnegative
right derivative. Conversely, the secant from zero to one lies above
that derivative, so the displayed conditions imply \(I(g)\geq I(f)\)
for every finite-information competitor. An infinite-information
competitor has a larger value already. This proves both directions.
\end{proof}
\end{annexitem}

\begin{annexitem}{annex:p15.information.saddlepoint.criterion}{thm:p15-information-saddlepoint-criterion}{Exactly what is needed beyond a minimizing density}
\begin{proof}
For the density-based domain, integration against the efficient score
and Cauchy--Schwarz imply \(K_{\rm d}(\eta,F_*)\leq I(F_*)\)
for every admissible \(\eta\). Equality at positive information holds
exactly for a nonzero multiple of \(-f_*'/f_*\) on its support.
For the weak-derivative domain the same fact requires its integration
identity, which we now verify. If \(F=f\,dx\) has finite information,
\(\eta\) is locally absolutely continuous, its derivative representative
\(\chi\) is bounded, and \(\int\eta^2f<\infty\), then
\[
 \int|\eta f'|\leq
       \left(\int\eta^2f\right)^{1/2}\sqrt{I(F)}<\infty.
\]
Also \(\eta f\in L^1\) and \((\eta f)'=\chi f+\eta f'\in L^1\).
The limits of \(\eta f\) at both infinities exist and must be zero.
Integration gives \(\int\chi f=-\int\eta f'\).
The same Cauchy--Schwarz argument therefore applies.
Test scores approximate the information value at \(F_*\), so the
statistician's supremum there is exactly \(\alpha\).
The two inequalities in the statement are precisely the two
saddlepoint inequalities; this proves the equivalence and efficiency.

If \(F\ll F_*\), its density and its derivative vanish almost everywhere
outside \(\{f_*>0\}\). The preceding first-variation theorem has
\(r_g=0\); it forces \(b_g<\infty\) and \(2a_g-b_g\geq\alpha\).
For the efficient score at \(F_*\), \(a=b=\alpha\).
Since \(K\geq2a-b\), this proves the density-class saddlepoint.

Finally suppose a real continuous \(\psi\) on the whole line extends
the efficient score. The identity \(f_*'=-\psi f_*\) holds almost
everywhere, including its zero set. On every bounded interval multiply
by the positive integrating factor \(\exp(\int_0^x\psi(t)\,dt)\).
Its product with \(f_*\) is absolutely continuous with derivative zero.
Thus \(f_*(x)=c\exp(-\int_0^x\psi)\) on the whole line for a common
constant \(c\). Positive information excludes the zero measure,
so \(c>0\). Every finite-information density is then dominated by
\(F_*\). This supplies, rather than assumes, the support implication.
\end{proof}
\end{annexitem}

\begin{annexitem}{annex:p15.information.closure.extension}{thm:p15-information-closure-extension}{The source's extension to singular limits}
\begin{proof}
The preceding result makes \(f_*\) positive everywhere and gives
\[
 2a_\psi(F)-b_\psi(F)\geq\alpha\qquad(F\in C_1).
\]
It gives finiteness of \(b\) there as well. The integration identity
with the specified weak derivative was proved above.
In either case in the statement, \(\chi\) is bounded. This is immediate
for \(C_0\). A finite upper semicontinuous function is bounded above
on a compact set; nonnegativity and vanishing at infinity finish
the other case.

For \(\chi\in C_0\), compact cutoffs and the mass bound show that
\(F\mapsto a_\psi(F)\) is vaguely continuous. For the second case,
write \(M=\sup\chi\) and define
\[
 \chi_m(x)=\sup_y(\chi(y)-m|x-y|)_+ .
\]
These functions are nonnegative, \(m\)-Lipschitz, vanish at infinity,
and decrease to \(\chi\). Indeed any positive maximizing sequence
can use only \(|x-y|\leq M/m\); upper semicontinuity gives the
pointwise limit, while the vanishing tail of \(\chi\) gives that
of \(\chi_m\). Consequently, when \(F_n\to F\) vaguely,
\[
 \limsup_n\int\chi\,dF_n
 \leq\int\chi_m\,dF\quad\hbox{for every }m,
\]
and monotone convergence as \(m\to\infty\) gives upper
semicontinuity of \(a_\psi\).
Since \(\psi^2\) is continuous and nonnegative, increasing compact
cutoffs prove lower semicontinuity of \(b_\psi\).

Choose \(F_n\in C_1\) tending vaguely to \(F\in C\).
The uniform bound on \(a_\psi(F_n)\) and
\(b_\psi(F_n)\leq2a_\psi(F_n)-\alpha\) bound their second moments.
Lower semicontinuity gives a finite \(b_\psi(F)\), and the two
semicontinuities yield
\[
 2a_\psi(F)-b_\psi(F)
 \geq\limsup_n(2a_\psi(F_n)-b_\psi(F_n))\geq\alpha .
\]
If \(b=0\), this inequality gives \(a\geq\alpha/2>0\), so the
payoff is infinity under its explicit convention. Otherwise
\(a^2/b\geq2a-b\). At \(F_*\) the payoff is \(\alpha\); at any
admissible opposing score its payoff there is at most \(\alpha\).
This proves the full saddlepoint.

For the statistical statement note first that subprobabilities
converging vaguely to a proper probability converge on bounded
continuous functions. Choose a compact cutoff \(0\leq h\leq1\)
whose integral under the limit is at least \(1-\delta\).
Its convergent integrals bound the masses of the approximants
outside its support by \(2\delta\) eventually. Approximate the
bounded function by its product with a larger cutoff; the tails
are controlled by its bound times \(2\delta\). Let \(\delta\to0\).
Thus boundedness of \(\psi\) gives \(b(F_n)\to E_F\psi^2\).
For the absolutely continuous \(F_n\), nondecreasing \(\psi\)
has \(a(F_n)\geq0\), regardless of modifications of \(\chi\)
on Lebesgue-null sets. The saddlepoint bound gives
\(a(F_n)^2\geq\alpha b(F_n)\).
Taking limits with the assumed \(\limsup a(F_n)\leq d_F\) proves
\(E_F\psi^2/d_F^2\leq1/\alpha\).
Centering, the actual population derivative and the remaining
convex-estimator hypotheses are exactly those required to apply
\cref{thm:p15-convex-location-normality}. They are not consequences
of vague compactness alone.
\end{proof}
\end{annexitem}

\begin{annexitem}{annex:p15.huber.general.density}{thm:p15-huber-general-density}{The full tangent-tail construction and normalization}
\begin{proof}
Because \(g\) is continuous, \(S=\{g>0\}\) is an open interval.
On it \(h\) is twice continuously differentiable and \(h'\) is
continuous and nondecreasing. Integrability of \(g\) forces a negative
slope somewhere and a positive slope somewhere: if \(h'\geq0\)
throughout, \(g\) cannot decay toward the left end; if that end is
finite, continuity and \(g=0\) there give the same contradiction.
The other case is reflected. At a finite endpoint of \(S\),
\(g\to0\), so \(h\to\infty\). The two slope signs imply linear
growth of \(h\) at infinite endpoints. Thus its sublevel sets
are compact and it attains a finite minimum at some \(x_m\).

The set \(|h'|\leq k\) is a nonempty interval with extended endpoints.
At any finite endpoint its slope is the asserted \(-k\) or \(k\)
by continuity; such an endpoint lies inside \(S\).
Minimizing \(h(y)+k|x-y|\) in \(y\), its slopes on the two sides
of \(x\) show that the minimizer is \(x\) in the middle interval,
the lower tangent point to its left, and the upper tangent point
to its right. If a slope has a plateau, its tangent is the same
at every point on that plateau. This proves the displayed formula.
There is no tail branch at an infinite endpoint.
Its integral is the central mass plus the finite tangent-tail
terms, and is finite. The function \(h_k\) is convex, continuously
differentiable and has derivative \(\psi_k\) equal to the clipped
slope, with the indicated constant continuation outside \(S\).
On every bounded interval this derivative is absolutely continuous:
it is the clipping of the \(C^1\) function \(h'\) in the middle
and constant on the tails, with matching values. Its derivative
is nonnegative, and its total integral is its increment between
its two limiting slopes, at most \(2k\).

Increasing \(k\) increases \(h_k\), so it decreases \(Z(k)\).
For continuity, fix \(x\) and \(k\) in a compact subinterval of
\((0,\infty)\). Any minimizing \(y\) lies in a common compact
sublevel set of \(h\), by comparison with \(x_m\).
On that set the objectives differ by at most a constant times
the difference of their \(k\)'s. Thus \(h_k(x)\) is continuous in \(k\).
The integrable function \(q_{k_0}\) dominates for \(k\geq k_0>0\);
dominated convergence proves continuity of \(Z\).
For \(x\in S\), \(h_k(x)=h(x)\) once \(k\geq|h'(x)|\).
For \(x\notin S\), compact sublevel sets and the distance penalty
imply \(h_k(x)\to\infty\). Hence dominated convergence gives
\(Z(k)\to\int g=1\) as \(k\to\infty\).
On the other hand
\[
 h_k(x)\leq h(x_m)+k|x-x_m|,\qquad
 Z(k)\geq2g(x_m)/k,
\]
so \(Z(0+)=\infty\).

If \(Z(k_2)>1\), at least one endpoint at \(k_2\) is finite;
otherwise \(q_{k_2}=g\). For \(0<k_1<k_2\), on that outer tail
the two envelopes have different exponential decay rates.
Their ordered values are therefore strictly different on a set
of positive length. This proves strict decrease above one.
The intermediate value theorem now gives the unique normalizing
\(k\) for \(0<\varepsilon<1\).
Since \(h_k\leq h\), \(q_k\geq g\). Thus
\[
 H_*(dx)=\frac{(1-\varepsilon)(q_k-g)}{\varepsilon}\,dx
\]
is a probability measure and
\(f_*=(1-\varepsilon)g+\varepsilon\,dH_*/dx\).
The matching tangent values and slopes prove that \(f_*\)
is positive and \(C^1\) on the whole line.
\end{proof}
\end{annexitem}

\begin{annexitem}{annex:p15.huber.general.minimax}{thm:p15-huber-general-minimax}{The general saddlepoint and its exact statistical interpretation}
\begin{proof}
Write \(a=1-\varepsilon\), \(\psi=\psi_k\), and
\(\chi=\psi'\geq0\) almost everywhere.
The density \(f_*\) is bounded, \(\chi\in L^1(dx)\), and
\(|\psi|\leq k\). Integration by parts, with zero boundary term
\(\psi f_*\) at both infinities, gives
\[
 I_*=\int\psi^2f_*=\int\chi f_*=a\int\chi g.
\]
It is positive since a nonzero integrable density cannot have
identically zero information. The excess density \(q_k-g\) vanishes
on the middle interval; on either nonempty excess tail
\(\psi^2=k^2\) and \(\chi=0\) almost everywhere.
This proves the second formula for \(I_*\).

Let \(F= aG+\varepsilon H\) be in the vague closure.
If its information is finite, write \(F=f\,dx\).
The density characterization bounds \(f\), so
\(\int\chi f<\infty\). Since \(f\geq ag\) and \(H\) has mass at most one,
\[
 A:=\int\chi f\geq I_*,
 \qquad B:=\int\psi^2f\leq a\int\psi^2g+\varepsilon k^2=I_* .
\]
Integration by parts is justified by bounded \(\psi\) and
\(f'\in L^1\), giving \(A=-\int\psi f'\).
Cauchy--Schwarz yields \(I(F)\geq A^2/B\geq I_*\);
\(B>0\) follows from \(A>0\).
Infinite-information members cannot improve this value.
The positive-everywhere minimizer is unique by the exact equality
theorem. This argument does not require a bounded representative
of \(\chi\), unlike the separate abstract extension theorem.

We next prove the statistical assertion on the fixed centered class.
Let \(\lambda_F(t)=E_F\psi(X-t)\), so \(\lambda_F(0)=0\).
It is continuous and nonincreasing, by bounded continuous
nondecreasing \(\psi\). The prototype density is bounded and uniformly
continuous: it is continuous with zero limits at its support ends
and at infinity. Since \(\chi\in L^1\), Fubini gives
\[
 \lambda_G(t)-\lambda_G(0)
       =-\int_0^t\int\chi(x-u)g(x)\,dx\,du .
\]
The inner convolution is continuous at zero, by uniform continuity
of \(g\) and integrability of \(\chi\), with value \(I_*/a\).
The contaminant's score is nonincreasing. Thus, as \(t\to0\),
\[
 \lambda_F(t)\leq-I_*t+o(t)\quad(t>0),\qquad
 \lambda_F(t)\geq-I_*t+o(|t|)\quad(t<0).
\]
These inequalities give strict opposite signs near zero and,
by monotonicity, at every nonzero \(t\).
The convex potential \(h_k\) is coercive, with negative and positive
limiting slopes. Its finite-sample objective therefore has a nonempty
bounded interval of minimizers, whose midpoint is translation
equivariant. The score sign implications and the strong law for the
bounded score give consistency, as in the convex-estimator proof
in \cref{thm:p15-convex-location-normality}.
For completeness the endpoint computation follows from the same
monotonicity. The continuous sample score \(U_n\) has strictly
positive and negative limits at the two infinities, respectively.
Its zero set is a nonempty bounded closed interval \([a_n,b_n]\).
On a bracket with opposite signs, bisect and keep the half containing
the boundary between \(U_n>0\) and \(U_n\leq0\) to find \(a_n\).
Use the boundary between \(U_n\geq0\) and \(U_n<0\) to find \(b_n\).
Each interval still contains its endpoint and its length halves,
which proves convergence to the two endpoints and their midpoint.

Here is the risk argument even without a derivative for \(H\).
Put \(B=E_F\psi^2\leq I_*\). In fact \(B>0\), since \(F\geq aG\)
and \(\psi\) is not zero \(G\)-almost everywhere.
For fixed \(z>0\), boundedness and continuity of \(\psi\) imply
that the centered row variables \(\psi(X_i-z/\sqrt n)\)
have variance tending to \(B\), and their Lindeberg condition holds
because the variables are uniformly bounded.
The triangular-array central limit theorem
\cref{thm:p15-clt-lindeberg-feller} therefore applies.
If the midpoint estimate \(T_n>z/\sqrt n\), the score sum at
\(z/\sqrt n\) is nonnegative; the preceding mean bound gives
\[
 \limsup_n P_F(\sqrt n T_n>z)
      \leq 1-\Phi(I_*z/\sqrt B)
      \leq1-\Phi(\sqrt{I_*}\,z).
\]
One can make the first limit passage explicit by replacing its
threshold \(I_*z+o(1)\) by \(I_*z-\delta\) and then letting
\(\delta\downarrow0\). The reflected score implication at
\(-z/\sqrt n\) gives the identical bound for the lower tail.
Therefore
\[
 \limsup_n P_F(\sqrt n|T_n|>z)
       \leq2(1-\Phi(\sqrt{I_*}\,z)).
\]
Using
\(E\min(W^2,C^2)=\int_0^C2zP(|W|>z)\,dz\),
bounded reverse Fatou gives the same upper bound on limiting
truncated risk as for \(N(0,1/I_*)\).
Letting \(C\to\infty\) gives \(1/I_*\).
This is pointwise in \(F\) before taking the supremum in the
stated risk order; no uniform finite-\(n\) exchange is made.

At \(F_*\), the same convolution argument with its bounded uniformly
continuous density gives \(\lambda_{F_*}'(0)=-I_*\);
its score variance is \(I_*\) and its score mean is zero by
\(-\int f_*'=0\). The convex normality theorem gives variance
\(1/I_*\) and equality in the limiting truncated risk.
For any translation-equivariant estimator sequence, the
integrated-score lower bound
\cref{thm:p15-location-integrated-score-bound}, applied to the
positive \(C^1\) density \(f_*\), gives
\[
 E_{F_*}\min(nT_n^2,C^2)
       \geq\frac1{I_*+4\pi^2/C^2}.
\]
Taking the exact ordered limits establishes the all-equivariant
minimax assertion. For any other \(F\) having a finite population
derivative, the displayed one-sided bounds force \(d_F\geq I_*\);
bounded scores give the remaining moment-continuity conditions
of the normality theorem, so \(B/d_F^2\leq1/I_*\).
If \(g\) is symmetric, \(h_k\) is even and \(\psi\) odd, proving
centering for symmetric \(H\).
At \(\varepsilon=0\) the tail regularization disappears. None of
the bounded-score arguments is thereby a proof of an unbounded-score
normal limit without its stated hypotheses.
\end{proof}
\end{annexitem}

\begin{annexitem}{annex:p15.huber.nonnormal.examples}{thm:p15-huber-nonnormal-examples}{Logistic and asymmetric Gumbel calculations}
\begin{proof}
For the logistic law its CDF is
\((1+\tanh(x/2))/2\). Thus \(s=\tanh(X/2)\) is uniform on \((-1,1)\),
and the endpoints \(s=\pm k\) give
\(t=\log((1+k)/(1-k))\). The central mass is \(k\),
and \(g(t)=(1-k^2)/4\). Hence
\[
 Z(k)=k+\frac{1-k^2}{2k}=\frac{1+k^2}{2k},
 \qquad \varepsilon=1-Z(k)^{-1}.
\]
On the middle interval
\(\psi'(x)=\frac12(1-s^2)\), so
\[
 E_G\psi'=\frac14\int_{-k}^k(1-s^2)\,ds
             =\frac k2-\frac{k^3}6,\qquad
 E_G\psi^2=k^2-\frac{2k^3}3 .
\]
At \(k=1/2\), these are \(11/48\) and \(1/6\).
Multiplication of the first by \(4/5\), or addition of
\((4/5)(1/6)+(1/5)(1/4)\), gives \(I_*=11/60\).
For \(t=\log(7/5)\),
\(\tanh(t/2)=1/6\), so three zero observations contribute
\(-3/6\). If \(M>\log(21/5)=t+\log3\), the fourth residual
is on the upper clipped tail and contributes \(1/2\).
Their sum is zero. The three central terms have strictly negative
derivative with respect to \(t\) near this root; monotonicity
then makes it the unique root and minimizer.

For the Gumbel law use \(y=e^{-x}\). Its CDF is
\(e^{-e^{-x}}\), its density transforms as
\(g(x)\,dx=-e^{-y}\,dy\), and its score is \(1-y\).
At \(k=1\) the lower tangent point has \(y=2\), while the upper
slope approaches one without reaching it at a finite point.
The central mass is \(1-e^{-2}\); the lower tail contribution
is \(g(-\log2)=2e^{-2}\). Thus \(Z=1+e^{-2}\).
In the unclipped region \(\psi'=y\), giving
\[
 E_G\psi'=\int_0^2 ye^{-y}\,dy=1-3e^{-2},
 \qquad I_*=\frac{1-3e^{-2}}{1+e^{-2}} .
\]
For the centering restriction compute
\[
 E_G\psi=\int_0^2(1-y)e^{-y}\,dy-\int_2^\infty e^{-y}\,dy
         =2e^{-2}-e^{-2}=e^{-2}.
\]
Since \(\varepsilon/(1-\varepsilon)=e^{-2}\),
\((1-\varepsilon)E_G\psi+\varepsilon E_H\psi=0\)
requires \(E_H\psi=-1\).
The score is at least \(-1\), attaining it exactly on
\((-\infty,-\log2]\), so this is equivalent to the stated support
restriction. For the sample \((-M,0,0)\) and \(t=-\log2\),
the first residual is \(-M+\log2<-\log2\) and contributes \(-1\).
Each other residual is \(\log2\) and contributes \(1-1/2=1/2\).
The sum is zero, and the two central terms again make it the
unique root.
\end{proof}
\end{annexitem}

\begin{annexitem}{annex:p15.information.game.failures}{thm:p15-information-game-failures}{Support, derivative and attainment failures}
\begin{proof}
Explicitly, \(f(x)=2\sin^2(\pi x)\) on \((0,1)\), zero otherwise,
and \(g(x)=f(x-2)\). Their densities are absolutely continuous;
direct integration gives
\(\int f'^2/f=8\pi^2\int_0^1\cos^2(\pi x)\,dx=4\pi^2\),
and the same for \(g\).
Disjoint supports make information of every mixture equal to this
value. On \((0,1)\) the efficient score is
\(-2\pi\cot(\pi x)\); use it there and the constant one on \((2,3)\).
For \(f\), the density-based numerator and denominator both equal
\(4\pi^2\). For \(g\), the numerator is
\(-\int g'=0\) and the denominator is one, so the payoff is zero.
All required absolute integrals are finite by Cauchy--Schwarz.
If instead the score on \((2,3)\) is
\(-2\pi\cot(\pi(x-2))\), both numerator and denominator equal
\(4\pi^2\) for every mixture. These scores are used in the
density-based domain; no continuous extension through their poles
is asserted.

For the step score and normal density, direct integration gives
\(-\int_0^\infty\phi'(x)\,dx=\phi(0)\).
Its weak derivative is the point mass at zero, not its ordinary
almost-everywhere derivative zero. It fails the absolutely continuous
score premise of the derivative-representative domain.
Finally choose \(\eta\in C_c^1\) equal to one near zero and let
\(\psi_n(x)=\eta(x)(x+1/n)\). At \(F=\delta_0\),
\(A_{\psi_n}=1\), \(B_{\psi_n}=1/n^2\), so \(K(\psi_n,F)=n^2\).
This proves infinite information using positive denominators.
\end{proof}
\end{annexitem}
